\documentclass[11pt]{article}
% Structure and the things that point at it: sectioning with automatic
% numbers, a table of contents, labels resolved after the whole source is
% read (\ref, \eqref, \autoref, \cref), footnotes collected at the end,
% citations against a thebibliography list, an appendix, and a section
% spliced in from another file with \input.
\usepackage{amsmath}
\usepackage{hyperref}
\usepackage{cleveref}
\newcommand{\SI}[2]{$#1\,\mathrm{#2}$}
\title{Structure, References and Citations}
\author{The UltraCanvas Team\thanks{Written to be read by
  \texttt{UltraCanvasLaTeXDocumentReader}, not by \TeX.}}
\date{September 2026}
\begin{document}
\maketitle
\tableofcontents

\begin{abstract}
A reference is a forward reference as often as not, so the importer reads
the whole source first, collecting every \texttt{\textbackslash label} with
the number of the thing it was attached to, and resolves the references
afterwards. An undefined label prints \texttt{??} and is reported as a
diagnostic rather than silently swallowed.
\end{abstract}

\section{Sectioning}\label{sec:sectioning}
Unstarred sectioning commands are numbered the way the \texttt{article}
class numbers them, and the numbers are part of the heading text.

\subsection{Depth}\label{sec:depth}
\texttt{\textbackslash section}, \texttt{\textbackslash subsection} and
\texttt{\textbackslash subsubsection} are headings 2, 3 and 4;
\texttt{\textbackslash part} and \texttt{\textbackslash chapter} are 1 above
them, and \texttt{\textbackslash paragraph} and
\texttt{\textbackslash subparagraph} are 5 and 6.

\subsubsection{Three levels down}
This heading carries the number \texttt{1.1.1}, and nothing but the counter
state produces it.

\paragraph{A run-in heading.} The paragraph command gives a heading that a
typesetter would set run-in with the text; here it is simply the deepest
heading level.

\section*{An unnumbered section}
A starred section takes no number and does not advance the counter, which is
why the next section is number~2.

% A fragment, not a document: no preamble and no \begin{document}. It is
% spliced into article-structure-refs.tex by % A fragment, not a document: no preamble and no \begin{document}. It is
% spliced into article-structure-refs.tex by % A fragment, not a document: no preamble and no \begin{document}. It is
% spliced into article-structure-refs.tex by \input{parts/units-note}, which
% resolves against the including file's directory. Fragments live in this
% subfolder so the demo's file scan, which does not descend into
% directories, lists only whole documents as tabs.
\subsection{A section from another file}\label{sec:units}
This subsection is written in \texttt{parts/units-note.tex} and spliced in
where the \texttt{\textbackslash input} stands, before anything else is
read. Numbering, labels and footnotes continue across the
splice\footnote{This footnote is defined in the included fragment and lands
in the same footnote list as the others.} --- the fragment cannot tell that
it is not part of the main file, which is the point of the mechanism.

The macro \verb|\SI| defined in the main preamble works here too:
\SI{299792458}{m/s} is the speed of light, and \SI{9.81}{m/s^2} the
standard gravity.
, which
% resolves against the including file's directory. Fragments live in this
% subfolder so the demo's file scan, which does not descend into
% directories, lists only whole documents as tabs.
\subsection{A section from another file}\label{sec:units}
This subsection is written in \texttt{parts/units-note.tex} and spliced in
where the \texttt{\textbackslash input} stands, before anything else is
read. Numbering, labels and footnotes continue across the
splice\footnote{This footnote is defined in the included fragment and lands
in the same footnote list as the others.} --- the fragment cannot tell that
it is not part of the main file, which is the point of the mechanism.

The macro \verb|\SI| defined in the main preamble works here too:
\SI{299792458}{m/s} is the speed of light, and \SI{9.81}{m/s^2} the
standard gravity.
, which
% resolves against the including file's directory. Fragments live in this
% subfolder so the demo's file scan, which does not descend into
% directories, lists only whole documents as tabs.
\subsection{A section from another file}\label{sec:units}
This subsection is written in \texttt{parts/units-note.tex} and spliced in
where the \texttt{\textbackslash input} stands, before anything else is
read. Numbering, labels and footnotes continue across the
splice\footnote{This footnote is defined in the included fragment and lands
in the same footnote list as the others.} --- the fragment cannot tell that
it is not part of the main file, which is the point of the mechanism.

The macro \verb|\SI| defined in the main preamble works here too:
\SI{299792458}{m/s} is the speed of light, and \SI{9.81}{m/s^2} the
standard gravity.


\section{Pointing at things}\label{sec:refs}
Section~\ref{sec:sectioning} is a plain \texttt{\textbackslash ref}. The
cleveref and hyperref spellings resolve to the same number and nothing more:
\autoref{sec:depth}, \Cref{sec:units} and \cref{eq:gauss} each give the
bare counter, so the word in front of it is the writer's, as in
Section~\ref{sec:depth}. Brackets around an equation number come from
\texttt{\textbackslash eqref}: \eqref{eq:gauss}.

\begin{equation}\label{eq:gauss}
  \int_{-\infty}^{\infty} e^{-x^2}\,dx = \sqrt{\pi}
\end{equation}

Numbers come from the counters the reader keeps as it goes: sections from
the sectioning commands, equations from the numbered math environments,
figures and tables from their captions, and items from the enumerate they
sit in.

\section{Notes and sources}
A footnote is a superscript number where it stands and an entry in a list
after a rule at the end of the document.\footnote{Like this one. Footnotes
keep their \emph{formatting}, and a footnote may hold a formula such as
$e^{i\pi} + 1 = 0$.} A second one\footnote{The numbering is the order in
which the notes are met, across included files as well.} follows here.

Citations are resolved against the bibliography below:
\cite{knuth84} is a single key, \cite{lamport94,mittelbach04} a list, and
\citep{knuth84} one of the natbib spellings. A key with no
\texttt{\textbackslash bibitem} keeps the key itself, so the source of the
gap is visible.

\appendix
\section{An appendix}\label{sec:appendix}
After \texttt{\textbackslash appendix} the section counter restarts and
counts in letters, so this is Appendix~\ref{sec:appendix}.

\begin{thebibliography}{9}
\bibitem{knuth84} D. E. Knuth, \emph{The \TeX book}, Addison-Wesley, 1984.
\bibitem{lamport94} L. Lamport, \emph{\LaTeX: A Document Preparation
  System}, 2nd ed., Addison-Wesley, 1994.
\bibitem{mittelbach04} F. Mittelbach and M. Goossens, \emph{The \LaTeX{}
  Companion}, 2nd ed., Addison-Wesley, 2004.
\end{thebibliography}
\end{document}
